\documentclass[12pt,a4paper]{article}

\usepackage[utf8]{inputenc}
\usepackage[T1]{fontenc}
\usepackage[brazil]{babel}
\usepackage{geometry}
\usepackage{booktabs}
\usepackage{array}
\usepackage{float}
\usepackage{graphicx}
\usepackage{hyperref}
\usepackage{tikz}

\geometry{margin=2.5cm}
\hypersetup{
    colorlinks=true,
    linkcolor=black,
    urlcolor=blue
}

\newcommand{\optionalfigure}[2]{%
\begin{figure}[H]
    \centering
    \IfFileExists{figuras/#1}{%
        \includegraphics[width=0.92\textwidth]{figuras/#1}%
    }{%
        \IfFileExists{relatorio/figuras/#1}{%
            \includegraphics[width=0.92\textwidth]{relatorio/figuras/#1}%
        }{%
            \fbox{\parbox{0.86\textwidth}{Imagem opcional: #2.}}%
        }%
    }%
    \caption{#2}
\end{figure}
}

\title{\textbf{Projeto QoE SDN -- Etapa 1}\\
\large Melhoria da QoE em Streaming de Vídeo com Mininet, SDN e P4}
\author{Disciplina de Programabilidade de Redes\\
Repositório: \texttt{ARTHUR9011/projeto-qoe-sdn}}
\date{Maio de 2026}

\begin{document}

\maketitle
\thispagestyle{empty}
\newpage

\tableofcontents
\newpage

\section{Introdução}

Este relatório descreve a Etapa 1 do projeto ``Melhoria da QoE em Streaming de Vídeo com Mininet, SDN e P4''. O foco desta entrega é montar um ambiente experimental funcional para streaming DASH em uma rede SDN emulada, com controlador OpenFlow, topologia Mininet e medições iniciais de conectividade e vazão.

\section{Objetivo da Etapa 1}

O objetivo da Etapa 1 é construir uma base reproduzível para experimentos futuros de QoE. A etapa inclui a configuração de uma topologia com um servidor de vídeo, três clientes, um switch OpenFlow 1.3 e um controlador SDN remoto. Também inclui a validação do acesso HTTP ao manifesto DASH e a coleta de métricas de rede com \texttt{ping} e \texttt{iperf3}.

\section{Fundamentação aplicada}

SDN separa o plano de controle do plano de dados. O controlador SDN executa a lógica de decisão e programa o comportamento dos switches, enquanto os switches encaminham pacotes conforme as regras instaladas.

O OpenFlow é usado como protocolo entre controlador e switch. Ele permite instalar regras em tabelas de fluxo, associando condições de casamento a ações de encaminhamento.

O Mininet permite emular hosts, switches, enlaces e controladores em Linux, usando namespaces de rede e Open vSwitch. Isso torna possível reproduzir uma rede SDN completa em uma única máquina.

P4 está relacionado à programabilidade do plano de dados e pode ser considerado em etapas futuras. Nesta Etapa 1, P4 é apenas citado como extensão possível e não foi implementado.

\section{Ambiente experimental}

A topologia implementada contém:

\begin{itemize}
    \item \texttt{h1}: servidor DASH, IP \texttt{10.0.0.1/24};
    \item \texttt{h2}: cliente 1, IP \texttt{10.0.0.2/24};
    \item \texttt{h3}: cliente 2, IP \texttt{10.0.0.3/24};
    \item \texttt{h4}: cliente 3, IP \texttt{10.0.0.4/24};
    \item \texttt{s1}: switch OpenFlow 1.3;
    \item \texttt{c0}: controlador SDN remoto em \texttt{127.0.0.1:6633}.
\end{itemize}

Todos os enlaces usam \texttt{TCLink} com banda inicial de 100 Mbps. Não foram configurados atraso, perda, jitter ou degradação artificial nesta etapa.

\section{Ferramentas utilizadas}

\begin{itemize}
    \item Linux;
    \item Mininet;
    \item Open vSwitch;
    \item Python 3;
    \item OS-Ken, com compatibilidade de código para Ryu;
    \item \texttt{iperf3};
    \item \texttt{curl};
    \item \texttt{ffmpeg}, caso seja necessário regenerar os segmentos DASH.
\end{itemize}

O VLC é uma opção complementar para validação visual, mas a validação registrada nesta entrega foi feita com \texttt{curl -I} no arquivo \texttt{manifest.mpd}.

\section{Topologia implementada}

\begin{figure}[H]
\centering
\begin{tikzpicture}[
    host/.style={draw, rounded corners, minimum width=2.2cm, minimum height=0.9cm},
    switch/.style={draw, rectangle, minimum width=2.2cm, minimum height=0.9cm},
    controller/.style={draw, dashed, rounded corners, minimum width=3.2cm, minimum height=0.9cm}
]
    \node[host] (h1) at (-4,0) {\texttt{h1}\\10.0.0.1};
    \node[switch] (s1) at (0,0) {\texttt{s1}\\OpenFlow 1.3};
    \node[host] (h2) at (0,2.4) {\texttt{h2}\\10.0.0.2};
    \node[host] (h3) at (4,0) {\texttt{h3}\\10.0.0.3};
    \node[host] (h4) at (0,-2.4) {\texttt{h4}\\10.0.0.4};
    \node[controller] (c0) at (0,4.2) {\texttt{c0}\\127.0.0.1:6633};

    \draw (h1) -- node[above]{100 Mbps} (s1);
    \draw (s1) -- node[right]{100 Mbps} (h2);
    \draw (s1) -- node[above]{100 Mbps} (h3);
    \draw (s1) -- node[right]{100 Mbps} (h4);
    \draw[dashed] (c0) -- node[right]{OpenFlow} (s1);
\end{tikzpicture}
\caption{Topologia da Etapa 1.}
\end{figure}

\section{Controlador SDN}

O controlador está em \texttt{controlador/simple\_switch\_13.py}. Ele implementa um switch L2 simples com aprendizado de MAC em OpenFlow 1.3:

\begin{itemize}
    \item aprende o MAC de origem de cada pacote recebido;
    \item encaminha para a porta conhecida quando o MAC de destino já foi aprendido;
    \item usa \texttt{FLOOD} quando o destino ainda é desconhecido;
    \item instala regras de fluxo para destinos conhecidos.
\end{itemize}

O comando principal de execução é:

\begin{verbatim}
osken-manager --ofp-tcp-listen-port 6633 controlador/simple_switch_13.py
\end{verbatim}

\section{Streaming DASH}

O vídeo de teste e os segmentos DASH ficam em \texttt{video/}. O manifesto está em \texttt{video/dash/manifest.mpd}, com segmentos \texttt{.m4s}. O servidor HTTP usado na Etapa 1 é o servidor embutido do Python 3, executado no host \texttt{h1}:

\begin{verbatim}
h1 bash -c "cd video/dash && python3 -m http.server 8000 --bind 10.0.0.1" &
\end{verbatim}

O acesso dos clientes é validado por:

\begin{verbatim}
h2 curl -I http://10.0.0.1:8000/manifest.mpd
h3 curl -I http://10.0.0.1:8000/manifest.mpd
h4 curl -I http://10.0.0.1:8000/manifest.mpd
\end{verbatim}

\section{Metodologia de validação}

A validação da Etapa 1 considera:

\begin{itemize}
    \item conectividade geral com \texttt{pingall};
    \item latência e perda com \texttt{ping} entre servidor e cliente;
    \item vazão com \texttt{iperf3};
    \item disponibilidade do streaming DASH com \texttt{curl -I} no manifesto.
\end{itemize}

O script \texttt{scripts/run\_baseline.sh} automatiza a coleta enquanto controlador e topologia estão ativos. Ele salva os arquivos de saída em \texttt{resultados/}.

\section{Resultados iniciais}

Os resultados registrados na pasta \texttt{resultados/} indicam funcionamento estável no cenário base:

\begin{table}[H]
\centering
\begin{tabular}{p{3cm} p{7.4cm} p{3.2cm}}
\toprule
Teste & Resultado observado & Arquivo \\
\midrule
\texttt{h1 -> h2} ping & 10 enviados, 10 recebidos, 0\% perda, RTT médio 0,123 ms & \texttt{ping\_h1\_h2.txt} \\
\texttt{h2 -> h1} ping & 10 enviados, 10 recebidos, 0\% perda, RTT médio 0,078 ms & \texttt{ping\_h2\_h1.txt} \\
\texttt{h2 -> h1} iperf3 & receiver aproximadamente 70,3 Mbits/sec & \texttt{iperf\_h2\_h1.txt} \\
\texttt{h3 -> h1} iperf3 & receiver aproximadamente 72,9 Mbits/sec & \texttt{iperf\_h3\_h1.txt} \\
\texttt{h4 -> h1} iperf3 & receiver aproximadamente 68,9 Mbits/sec & \texttt{iperf\_h4\_h1.txt} \\
DASH via \texttt{h2}, \texttt{h3}, \texttt{h4} & HTTP 200 OK, \texttt{application/dash+xml}, 2201 bytes & \texttt{dash\_h*\_header.txt} \\
\bottomrule
\end{tabular}
\caption{Medições iniciais sem degradação artificial.}
\end{table}

\section{Evidências dos testes}

As evidências textuais estão na pasta \texttt{resultados/}. A validação DASH foi feita com \texttt{curl -I} no manifesto e não há evidência registrada de uso do VLC nesta etapa.

\optionalfigure{estrutura_repositorio.png}{Print da estrutura do repositório}
\optionalfigure{controlador_rodando.png}{Print do controlador rodando}
\optionalfigure{topologia_nodes_net_dump.png}{Print da topologia Mininet com \texttt{nodes}, \texttt{net} e \texttt{dump}}
\optionalfigure{pingall.png}{Print do \texttt{pingall}}
\optionalfigure{dash_curl.png}{Print do acesso DASH com \texttt{h2}, \texttt{h3} e \texttt{h4}}
\optionalfigure{iperf3.png}{Print do \texttt{iperf3} ou dos arquivos de resultado}

\section{Organização do repositório}

\begin{itemize}
    \item \texttt{controlador/}: controlador SDN da Etapa 1;
    \item \texttt{topologia/}: topologia Mininet;
    \item \texttt{scripts/}: scripts de execução, baseline e limpeza;
    \item \texttt{video/}: vídeo de teste e arquivos DASH;
    \item \texttt{resultados/}: saídas de ping, iperf3 e validação HTTP;
    \item \texttt{relatorio/}: relatório LaTeX.
\end{itemize}

\section{Instruções de reprodução}

\begin{enumerate}
    \item Instalar dependências: Mininet, Open vSwitch, Python 3, OS-Ken, \texttt{iperf3} e \texttt{curl}.
    \item Executar \texttt{make permissions}.
    \item Em um terminal, executar \texttt{make controller}.
    \item Em outro terminal, executar \texttt{make topology}.
    \item No CLI do Mininet, iniciar o servidor DASH com \texttt{h1 bash scripts/start\_dash\_server.sh \&}.
    \item Validar conectividade com \texttt{pingall}.
    \item Validar o DASH com \texttt{curl -I} a partir de \texttt{h2}, \texttt{h3} e \texttt{h4}.
    \item Validar vazão com \texttt{iperf3}.
    \item Opcionalmente, executar \texttt{make baseline} em um terceiro terminal para salvar os resultados automaticamente.
\end{enumerate}

\section{Verificação dos requisitos}

\begin{table}[H]
\centering
\begin{tabular}{p{4.2cm} p{7cm} p{2.4cm}}
\toprule
Requisito & Implementação no projeto & Situação \\
\midrule
Mininet em Linux & Topologia em \texttt{topologia/topologia\_qoe.py} com hosts, switch e enlaces \texttt{TCLink}. & Atendido \\
Controlador SDN & Controlador L2 em \texttt{controlador/simple\_switch\_13.py}, usando OpenFlow 1.3. & Atendido \\
Topologia com 1 servidor, 3 clientes e 1 switch OpenFlow & \texttt{h1} como servidor, \texttt{h2-h4} como clientes e \texttt{s1} como switch. & Atendido \\
Streaming DASH & Servidor Python no \texttt{h1} servindo \texttt{video/dash/manifest.mpd}. & Atendido \\
Métricas com ping e iperf3 & Arquivos em \texttt{resultados/} e automação em \texttt{scripts/run\_baseline.sh}. & Atendido \\
Automação com scripts/Makefile & Alvos \texttt{controller}, \texttt{topology}, \texttt{permissions}, \texttt{clean} e \texttt{baseline}. & Atendido \\
Relatório com diagrama e instruções & Este relatório inclui diagrama TikZ, metodologia e reprodução. & Atendido \\
\bottomrule
\end{tabular}
\caption{Checklist dos requisitos da Etapa 1.}
\end{table}

\section{Conclusão}

A Etapa 1 entrega uma base funcional para experimentos de streaming DASH sobre SDN em Mininet. O projeto inclui controlador SDN com OpenFlow 1.3, topologia reproduzível, servidor DASH, medições iniciais e automação por scripts e Makefile. As próximas etapas podem introduzir degradação controlada, tráfego concorrente, análise de QoE e, futuramente, programabilidade do plano de dados com P4.

\end{document}
